\section{Status i plan}\label{sec:status}

Stan na commitach z sekcji~\ref{sec:wprowadzenie-wersja}, według planów pracy
repozytoriów. Status zmienia się z każdym scalonym pull requestem; aktualny jest w
\path{docs/WORKPLAN.md} danego repozytorium.

\subsection{Faza 01}\label{sec:status-faza01}

\begin{xltabular}{\linewidth}{@{}>{\raggedright\arraybackslash}p{3.6cm} >{\raggedright\arraybackslash}X@{}}
\toprule
Repozytorium & Stan ticketów fazy 01 \\
\midrule
\endfirsthead
\toprule
Repozytorium & Stan ticketów fazy 01 \\
\midrule
\endhead
rdzeń & S0--S15 zrobione; S16 (outbound stubs) przeniesione do fazy 02, bo ani sklep
przykładowy, ani demonstracyjna konfiguracja nopCommerce nie wykonują wywołań
wychodzących. Nic nie jest wydane. \\
Portcullis & R0 (import), R1 (zmiana nazwy, pilnowana w CI), R2 (reguły migracyjne z
testami mutacyjnymi), R3 (SARIF z filtrem zmienionych linii i baseline) zrobione;
wsparcie P8 zaplanowane. Nic nie jest opublikowane: nie ma tagu \repopath{v*}, pakietów
w~nuget.org ani opublikowanego kontenera, a~workflow wydania nie miał jeszcze przebiegu;
publikacja przy pierwszym tagu \repopath{v*}. Według README pakiety i~obraz kontenera są
zbudowane i~sprawdzone lokalnie, ale CI buduje tylko pakiet analyzerów --- obrazu nie. \\
standardy & C4-a (standardy w Markdown z metadanymi), C4-b (generator z trybem
\repopath{--check}), C4-c (drift check) zrobione; R4 (skill dla modernize-dotnet)
scalony, zrobiony wtedy, gdy P6 potwierdzi, że agent go podjął. Żadna wersja nie jest
jeszcze otagowana, a~odnośniki do źródeł w~stopce skilla wskazują tag \texttt{v0.1.0},
którego nie ma. Sprawdzenie wykonywalne ma 11 z~20 standardów; pozostałe 9 to instrukcje
dla agenta. \\
bench nopCommerce & P1--P5 zrobione. P9 (bramka fazy 01) w toku: zabite mutanty
zrobione (\#8, 11 z 11); pakiet dowodowy kandydata, różnice i regresje czekają na
P6--P8. P10 (log godzin) w toku: godziny agenta do P5 i dla mutantów P9 są w logu, część
oszacowana po fakcie; godzin człowieka jeszcze nie ma. P6 (praca człowieka z GitHub
Copilot), P7 i P8 zaplanowane. \\
legacy-risk-scan & L18 scalony (\#1, \#2) i~zrobiony: strona działa na GitHub Pages od
26~IX~2026 (\url{https://konradcinkusz.github.io/letsgolegacy.legacy-risk-scan/}), a~job
\repopath{verify} w~\path{deploy.yml} przechodzi --- na \texttt{4f4e346} w~przebiegu
37116401196 strona serwuje ten commit, a~test z~prawdziwym OSV przechodzi
(podrozdział~\ref{sec:risk-scan-deploy}). Plan pracy odnotowuje wdrożenie od
\texttt{6bb73b1} (PR~\#9): wiersz L18 ma status „done”. Otwarte pozostają dwa kroki z~planu: ustawienie pozwalające
Actions otwierać pull requesty (miesięczne odświeżanie danych EOL z~1~X~2026 wypchnęło
gałąź, ale nie otworzyło PR) i~kontakt na stronie (ticket L17; dziś „Kontakt: wkrótce”). \\
capture, sandbox, portal & tylko plany; decyzja o topologii demonstracyjnej (D7) w stanie
„proposed”, zrobiona, gdy ADR zostanie przyjęty. \\
\bottomrule
\end{xltabular}

\zrodla{\path{letsgolegacy.secondkey/docs/WORKPLAN.md},
\path{letsgolegacy.secondkey/README.md},
\path{letsgolegacy.portcullis/docs/WORKPLAN.md},
\path{letsgolegacy.portcullis/docs/TUTORIAL.md},
\path{letsgolegacy.portcullis/README.md},
\path{letsgolegacy.portcullis/CHANGELOG.md},
\path{letsgolegacy.secondkey-standards/docs/WORKPLAN.md},
\path{letsgolegacy.secondkey-standards/README.md},
\path{letsgolegacy.secondkey-nopcommerce-bench/docs/WORKPLAN.md},
\path{letsgolegacy.legacy-risk-scan/docs/WORKPLAN.md},
\path{letsgolegacy.legacy-risk-scan/.github/workflows/deploy.yml},
\path{letsgolegacy.legacy-risk-scan/.github/workflows/update-eol.yml},
\path{letsgolegacy.secondkey-capture/docs/WORKPLAN.md},
\path{letsgolegacy.secondkey-sandbox/docs/WORKPLAN.md},
\path{letsgolegacy.secondkey-portal/docs/WORKPLAN.md}}

\subsection{Liczby, które już są}\label{sec:status-liczby}

\begin{xltabular}{\linewidth}{@{}>{\raggedright\arraybackslash}p{4.2cm} >{\raggedright\arraybackslash}X@{}}
\toprule
Co & Wynik zapisany przez CI \\
\midrule
\endfirsthead
\toprule
Co & Wynik zapisany przez CI \\
\midrule
\endhead
P3 --- rozgrzewka na eShopLegacyMVC & pakiet dowodowy wyprodukowany przez cały łańcuch (capture, A/A replay, werdykt, bramka, pakiet): 7 scenariuszy, 27 wymian; A/A 54 wyniki, 14 resetów, 0 bez odpowiedzi; werdykt \repopath{pass} --- 27 \repopath{equal}, 0 regresji; kontrakt 13 klauzul, wszystkie zaakceptowane i utrzymane, 0 niesprawdzonych, udział nieobecności 23\,\% (3 \repopath{never}); bramka na całym drzewie legacy: 27 znalezisk, \repopath{sk gate} \repopath{fail} z 7 blokującymi --- zapisane, nie egzekwowane \\
P4 --- zestaw ruchu nopCommerce & 38 scenariuszy (30 zaplanowanych, niektóre jako kilka sesji), 290 wymian \\
P5 --- kontrakt i dowód A/A & werdykt \repopath{pass}: 76 z 76 zaakceptowanych klauzul utrzymanych, 0 niesprawdzonych; 15 klauzul \repopath{never} (udział nieobecności 19,7\,\%); 290 wymian w 38 scenariuszach: 272 \repopath{equal}, 18 \repopath{equal-under-contract}, 0 regresji, 0 fix candidates \\
P9 --- zabite mutanty & 11 z 11: każdy z M01--M11 daje w \repopath{sk compare} wynik \repopath{fail}, a każda klauzula zarejestrowana dla mutanta z góry na nim zregresowała; M00 (kalibracja: źródło bez zmian) \repopath{pass} --- 290 wymian, 269 \repopath{equal}, 21 \repopath{equal-under-contract}, 0 regresji; to wynik dla tych jedenastu defektów na tym ruchu, a nie mutation score --- zasięg kontraktu zmierzą generowane mutanty \repopath{sk mutate} (faza 02) \\
Testy mutacyjne rdzenia & Stryker.NET w~CI na \texttt{c3ae5c2} (przebieg 37116794634), pięć z~siedmiu projektów produktu: \repopath{SecondKey.Compare} 98,89\,\%, \repopath{SecondKey.Evidence} 91,67\,\%, \repopath{SecondKey.Replay} (bez kodu SQL Server) 89,71\,\%, \repopath{SecondKey.Artifacts} 89,09\,\%, \repopath{SecondKey.Contract} 86,87\,\%; progi 85, 70 i~60 (high, low, break), poniżej 60 CI zawodzi; poza zakresem są \repopath{SecondKey.Capture} i~\repopath{SecondKey.Cli}. Wyniki są w~logach i~artefaktach CI, nie w~repozytorium; wynik \repopath{SecondKey.Replay} wzrósł z~60,7\,\% do 89,5\,\% w~drugiej rewizji analizy fazy 01 \\
Testy mutacyjne reguł migracyjnych Portcullis & 100,00\,\% w ostatnim zapisanym przebiegu Stryker.NET: 214 z 214 testowanych mutantów wykrytych (212 zabitych, 2 przekroczenia czasu); \path{mutation.yml} zawodzi poniżej progu 95 \\
\bottomrule
\end{xltabular}

Liczby P3--P5 i~P9 pochodzą z~przebiegów na starych pinach narzędzi (\texttt{967c49c},
\texttt{84e1925}; sekcja~\ref{sec:wprowadzenie-wersja}). Po przesunięciu pinów (ADR 0008)
workflowy benchu przechodzą na zielono, ale nowych liczb szczegółowych nie zapisano
w~\path{results/}.

\zrodla{\path{letsgolegacy.secondkey-nopcommerce-bench/docs/WORKPLAN.md},
\path{letsgolegacy.secondkey-nopcommerce-bench/results/p3/README.md},
\path{letsgolegacy.secondkey-nopcommerce-bench/results/p4/README.md},
\path{letsgolegacy.secondkey-nopcommerce-bench/results/p5/README.md},
\path{letsgolegacy.secondkey-nopcommerce-bench/results/p9/README.md},
\path{letsgolegacy.secondkey/docs/analysis/phase-01.md},
\path{letsgolegacy.secondkey/stryker-config.json},
\path{letsgolegacy.secondkey/scripts/mutation.sh},
\path{letsgolegacy.portcullis/docs/MUTATIONS.md},
\path{letsgolegacy.portcullis/tests/Portcullis.Engine.Tests/stryker-config.json}}

\subsection{Droga do bramki fazy 01}\label{sec:status-droga}

Pozostała ścieżka krytyczna przechodzi przez bench: P6 (migracja przez agenta, praca
człowieka) → P7 (replay legacy kontra kandydat; \repopath{verdict.json} istnieje, a różnica
NLS → ICU jest znaleziona albo jej brak udokumentowany) → P8 (reguły Portcullis na PR
kandydata; SARIF w pakiecie) → P9 (publikacja pakietu, kontraktu i liczb: scenariusze,
różnice, regresje, zabite mutanty). Zabite mutanty już są (podrozdział~\ref{sec:bench-p9}),
więc z liczb P9 brakuje już tylko tych, które przynosi kandydat. P10 zbiera godziny wdrożenia dla P4--P7; godziny
człowieka przyjdą z P6. Runbook P6 jest napisany przed pierwszym uruchomieniem; zapis
pierwszego przebiegu stanie się jego przykładem.

\zrodla{\path{letsgolegacy.secondkey-nopcommerce-bench/docs/WORKPLAN.md},
\path{letsgolegacy.secondkey-nopcommerce-bench/docs/P6-RUNBOOK.md},
\path{letsgolegacy.secondkey/docs/WORKPLAN.md}}

\subsection{Później}\label{sec:status-pozniej}

Zakres faz 02--04 jest w sekcji~\ref{sec:architektura-fazy}; szczegóły repozytoriów fazy
02 i 03 w rozdziale~\ref{sec:faza02}. Wszystkie te prace zaczynają się dopiero po bramce
fazy 01.

\zrodla{\path{letsgolegacy.secondkey/docs/WORKPLAN.md}}
